\section*{अध्याय \ref{ch:b3:gfd-foundations} --- भूभौतिकीय द्रव गतिकी}

\begin{solution}{exo:b3:gfd-foundations:1}
(क)~$f = 2\Omega\sin\varphi$: \num{7.292e-5}, \num{1.031e-4} और
\qty{1.263e-4}{s^{-1}}; $2\pi/f$: 23.93, 16.92 और \qty{13.82}{h}।
(ख)~$U/f$: 2.74, 1.94 और \qty{1.58}{km}। (ग)~\ang{30} पर, $\sin\varphi =
1/2$ और $2\pi/f = 2\pi/\Omega$: जड़त्वीय आवर्तकाल नक्षत्र दिवस,
\qty{23.93}{h} होता है।
\end{solution}

\begin{solution}{exo:b3:gfd-foundations:2}
$f = \qty{1.031e-4}{s^{-1}}$ मानने पर: धूल-भँवर \num{9.7e3}, समुद्र समीर 2.4,
अवदाब 0.15, वलय \num{2.4e-4}। अवदाब और वलय भू-विक्षेपी
संतुलन के निकट हैं; धूल-भँवर पृथ्वी के घूर्णन से बेखबर,
दोनों में से किसी भी दिशा में घूम सकता है, और समुद्र समीर इसे केवल एक दिन के दौरान महसूस करती है।
\end{solution}

\begin{solution}{exo:b3:gfd-foundations:3}
$\partial u_g/\partial z = (g/fT)\,|\partial T/\partial y| = 9.81\times
\num{5e-6}/(\num{1.031e-4}\times250) = \qty{1.90e-3}{s^{-1}}$; ऊँचाई में
\qty{10}{km} पश्चिमी पवन में \qty{19}{m/s} की वृद्धि होती है।
\end{solution}

\begin{solution}{exo:b3:gfd-foundations:4}
$E = \num{50e-8}$: $N = \qty{2.21e-3}{s^{-1}}$, आवर्तकाल
\qty{47}{min}, \num{1.27} चक्र प्रति घंटा। $E = \num{1000e-8}$: $N = \qty{9.90e-3}{s^{-1}}$,
\qty{10.6}{min}, 5.7 चक्र प्रति घंटा। $E = \num{e-8}$: $N =
\qty{3.13e-4}{s^{-1}}$, \qty{5.6}{h}, 0.18 चक्र प्रति घंटा।
\end{solution}

\begin{solution}{exo:b3:gfd-foundations:5}
(क)~$W \sim UH/L = \qty{0.1}{m/s}$ और $\mathrm Dw/\mathrm Dt \sim UW/L =
\qty{e-6}{m/s^2}$, $g$ का \num{1.0e-7} गुना। दाब पद, $fUL/H =
\qty{0.103}{m/s^2}$, \num{1e5} गुना बड़ा है: अनुपात, \num{9.7e-6},
$\mathrm{Ro}(H/L)^2$ है, जहाँ $\mathrm{Ro} = 0.097$। (ख)~$\tilde fu =
2\Omega\cos\ang{45}\times10 = \qty{1.03e-3}{m/s^2}$, $g$ का \num{1.05e-4} गुना।
\end{solution}

\begin{solution}{exo:b3:gfd-foundations:6}
$fr = \qty{30}{m/s}$। (क)~$\mathrm{Ro}_g = 0.667$, $V = 40/(1 + \sqrt{3.667})
= \qty{13.7}{m/s}$, भू-विक्षेपी पवन की तुलना में एक-तिहाई धीमी। (ख)~$\mathrm{Ro}_g =
0.167$, $V = 10/(1 + \sqrt{0.333}) = \qty{6.34}{m/s}$। (ग)~$V_g \le fr/4 =
\qty{7.5}{m/s}$, और तब $V = fr/2 = \qty{15}{m/s}$, भू-विक्षेपी
गति का दोगुना।
\end{solution}

\begin{solution}{exo:b3:gfd-foundations:7}
(क)~$\Omega = 2\pi/6 = \qty{1.047}{rad/s}$, $f = \qty{2.094}{s^{-1}}$।
$\mathrm{Ro} = 0.01/(2.094\times0.05) = 0.095$ और $\mathrm{Ek} =
\num{e-6}/(2.094\times0.04) = \num{1.19e-5}$: दोनों छोटे हैं, इसलिए प्रवाह
\omterm{thm:b3:gfd-foundations:taylor-proudman}{टेलर--प्राउडमैन प्रमेय} का पालन करता है और \omterm{thm:b3:gfd-foundations:taylor-proudman}{टेलर स्तंभ} बनता है। (ख)~$0.2/\sqrt{\num{e-6}
\times1.047} = \qty{195}{s}$, \qty{3.3}{min}, घूर्णन-मेज़ के 33 चक्कर।
\end{solution}

\begin{solution}{exo:b3:gfd-foundations:8}
$\partial\bar T/\partial y = -f\,\Delta u/(R\ln(850/250))$, जहाँ
$\ln(850/250) = 1.224$। \ang{30} N पर, $-\num{7.29e-5}\times36.12/(286.9
\times1.224) = \qty{-7.50}{K}$ प्रति \qty{1000}{km}; \ang{45} N पर,
$-\num{1.031e-4}\times17.18/351.1 = \qty{-5.05}{K}$ प्रति \qty{1000}{km}।
अपरूपण प्रवणता को $f$ से भाग देने पर समानुपाती होता है: प्रवणता 1.5
गुना बड़ी है \ang{30} पर, और $f$ 1.4 गुना छोटा है, इसलिए अपरूपण 2.1 गुना
बड़ा है।
\end{solution}

\begin{solution}{exo:b3:gfd-foundations:9}
(क)~$L_R = 0.0113\times11\,000/\num{1.031e-4} = \qty{1206}{km}$;
$0.00357\times1000/\num{9.47e-5} = \qty{37.7}{km}$, 32 गुना छोटा।
(ख)~$\beta = 2\Omega\cos\varphi/a$: \num{1.62e-11} और
\qty{1.74e-11}{m^{-1}.s^{-1}}। $\beta L_R/f = 0.19$ वायुमंडल के लिए, 0.007
महासागर के लिए।
\end{solution}

\begin{solution}{exo:b3:gfd-foundations:10}
(क)~$\mathrm dw/\mathrm dt = -(r + \mathrm if)w$, अतः $w = U\mathrm
e^{-rt}\mathrm e^{-\mathrm ift}$: वेग दर $f$ पर घूमता है, जबकि
इसका परिमाण $\mathrm e^{-rt}$ के रूप में क्षीण होता है। स्थिति $x + \mathrm iy =
U(1 - \mathrm e^{-(r + \mathrm if)t})/(r + \mathrm if)$ कुंडलित होकर
बिंदु $U/(r + \mathrm if)$ की ओर जाती है। (ख)~$x + \mathrm iy = U(1 - \mathrm e^{-(r +
\mathrm if)t})/(r + \mathrm if) + u_0t$: कुंडली लूपों में खिंच जाती है।
(ग)~$0.2\,\mathrm e^{-1} = \qty{0.074}{m/s}$; केंद्र $u_0t =
0.05\times4\times86\,400 = \qty{17.3}{km}$ पूर्व की ओर खिसक गया है।
\end{solution}

\begin{solution}{exo:b3:gfd-foundations:11}
$f = \qty{9.37e-5}{s^{-1}}$; $\partial u/\partial z = (g/f\rho_0)\,\partial
\rho/\partial y = 9.81\times10^{-5}/(\num{9.37e-5}\times1025) =
\qty{1.02e-3}{s^{-1}}$। \qty{1000}{m} पर $u = 0$ मानने पर सतही धारा
\qty{1.02}{m/s} है, पूर्व की ओर: उत्तरी गोलार्ध में धारा
हल्के जल को अपनी दाईं ओर रखकर बहती है।
\end{solution}

\begin{solution}{exo:b3:gfd-foundations:12}
(क)~$N^2 = (g/\rho)\,\Delta\sigma_\theta/\Delta z$। 0--\qty{100}{m}:
$9.81\times0.668/(1026.4\times100)$, $N = \qty{7.99e-3}{s^{-1}}$, 4.6
चक्र प्रति घंटा; 100--\qty{500}{m}: \qty{3.02e-3}{s^{-1}}, 1.7 चक्र
प्रति घंटा; 500--\qty{1000}{m}: \qty{3.24e-3}{s^{-1}}, 1.9 चक्र प्रति घंटा।
(ख)~ऊपरी \qty{100}{m} में वह गर्म सतही जल है जो मौसमी
घनत्व-प्रवणता परत के ऊपर है। 100 और \qty{500}{m} के बीच का जल
नीचे की तुलना में अधिक एकसमान है। (ग)~$E = N^2/g$ का मान \num{650e-8}, \num{93e-8} और
\qty{107e-8}{m^{-1}} है, सभी ऊपरी महासागर के परास के भीतर।
\end{solution}

\begin{solution}{pb:b3:gfd-foundations:1}
\textbf{1.} $f = 2\times\num{7.292115e-5}\sin\ang{48} =
\qty{1.084e-4}{s^{-1}}$; $2\pi/f = \qty{16.1}{h}$.

\textbf{2.} बाथटब \num{3.1e3}; टैंक ($f = 4\pi/10 = \qty{1.257}{s^{-1}}$)
0.020; समीर 2.3; अवदाब 0.14; वलय \num{2.3e-4}।

\textbf{3.} बाथटब को पृथ्वी के घूर्णन का कोई आभास नहीं होता; समीर
को इसका आभास केवल कुछ घंटों बाद होता है। टैंक, अवदाब और वलय की
रॉस्बी संख्याएँ छोटी हैं और भू-विक्षेपी संतुलन के निकट हैं।

\textbf{4.} टैंक $\num{e-6}/(1.257\times0.04) = \num{2.0e-5}$; अवदाब
$10/(\num{1.084e-4}\times10^8) = \num{9.2e-4}$; वलय
$\num{e-3}/(\num{1.084e-4}\times10^6) = \num{9.2e-6}$।

\textbf{5.} $\delta/H = \sqrt{2\nu/(fH^2)} = \sqrt{2\,\mathrm{Ek}}$। टैंक
$\sqrt{2\times\num{e-6}/1.257} = \qty{1.3}{mm}$; वायु $\sqrt{20/\num{1.084e-4}}
= \qty{430}{m}$।

\textbf{6.} टैंक और महासागर में गहराई के कुछ हज़ारवें भाग की परतों में
और वायुमंडल में कुछ सौवें भाग की परतों में घर्षण कार्य करता है; इनसे ऊपर
प्रवाह लगभग घर्षणरहित होता है।

\textbf{7.} $\rho = p/(RT) = 101\,290/(286.9\times288.1) =
\qty{1.225}{kg/m^3}$।

\textbf{8.} $u_g = (1/\rho f)\,|\partial p/\partial y| = (400/300\,000)/(1.225
\times\num{1.084e-4}) = \qty{10.0}{m/s}$, पश्चिम से, निम्न दाब
उत्तर में बाईं ओर होने पर।

\textbf{9.} $\mathrm{Ro}_g = 10.04/(\num{1.084e-4}\times600\,000) = 0.154$;
$V = 2\times10.04/(1 + \sqrt{1.618}) = \qty{8.84}{m/s}$, \qty{12}{\percent}
धीमी।

\textbf{10.} $V = 2\times10.04/(1 + \sqrt{0.382}) = \qty{12.4}{m/s}$।
सबसे बड़ा $V_g$ है $fr/4 = \qty{16.3}{m/s}$, एक प्रवणता $\rho fV_g =
\qty{2.16e-3}{Pa/m}$, \qty{6.5}{hPa} प्रति \qty{300}{km}।

\textbf{11.} $g/c_p = \qty{9.77}{K/km}$। $N^2 = (9.81/255.7)\times(9.77 -
6.49)\times10^{-3} = \qty{1.26e-4}{s^{-2}}$, $N = \qty{0.0112}{s^{-1}}$, जिसकी
अवधि \qty{9.3}{min} है।

\textbf{12.} $\Gamma = 0$ के साथ, $N^2 = (g/T)(g/c_p) = g^2/(c_pT) =
\qty{4.43e-4}{s^{-2}}$, $N = \qty{0.0210}{s^{-1}}$, जो
क्षोभमंडलीय मान का 1.9 गुना है।

\textbf{13.} $L_R = 0.0112\times10^4/\num{1.084e-4} = \qty{1035}{km}$, जो
अवदाब का आकार है।

\textbf{14.} $g' = \qty{0.0147}{m/s^2}$; $\sqrt{g'h} = \qty{0.54}{m/s}$;
$L_R = 0.542/\num{1.084e-4} = \qty{5.0}{km}$।

\textbf{15.} हाँ: झील, जो \qty{12}{km} चौड़ी है, अपनी
रॉस्बी त्रिज्या से दोगुने से अधिक है, इसलिए इसकी आंतरिक तरंगें और धाराएँ
कोरिऑलिस बल से मुड़ती हैं।

\textbf{16.} $(249.4 - 253.0)/600 = \qty{-6.0}{K}$ प्रति \qty{1000}{km}।

\textbf{17.} $\ln(1013/265) = 1.3409$। दक्षिण में:
$286.9\times253.0\times1.3409/9.81 = \qty{9922}{m}$; उत्तर में \qty{9781}{m}।
\qty{265}{hPa} सतह उत्तर में \qty{141}{m} नीचे है, निचली सतह के सापेक्ष
\num{2.35e-4} की ढलान।

\textbf{18.} $\Delta u = (g/f)\times\text{ढलान} = 9.81\times\num{2.35e-4}/
\num{1.084e-4} = \qty{21.3}{m/s}$।

\textbf{19.} $(R/f)\ln(p_b/p_t)\,|\partial\bar T/\partial y| =
(286.9/\num{1.084e-4})\times1.3409\times\num{6.0e-6} = \qty{21.3}{m/s}$, वही
है: दोनों परिकलन एक ही हैं।

\textbf{20.} \omterm{def:b3:gfd-foundations:baroclinic}{बैरोक्लिनिक}: तापमान, और इसलिए घनत्व, दाब-पृष्ठों के
साथ बदलता है। \omterm{def:b3:gfd-foundations:baroclinic}{बैरोट्रॉपिक} होने पर पवन ऊँचाई के साथ नहीं बदलती:
\qty{10.0}{m/s} होती।

\textbf{21.} $(g/f\bar T)\,|\partial T/\partial y|\,\Delta z = 9.81\times
\num{6.0e-6}\times10^4/(\num{1.084e-4}\times255.7) = \qty{21.2}{m/s}$,
\qty{0.3}{\percent} के भीतर।

\textbf{22.} $10.04 + (286.9/\num{1.084e-4})\times\num{6.0e-6}\times
\ln(1013/541) = 10.04 + 9.96 = \qty{20.0}{m/s}$।

\textbf{23.} $\partial u/\partial z = 21.3/10^4 = \qty{2.13e-3}{s^{-1}}$;
$\mathrm{Ri} = \num{1.26e-4}/\num{4.54e-6} = 28$, 0.25 से कहीं ऊपर: अपरूपण
मुक्त क्षोभमंडल को प्रक्षुब्ध नहीं करता।

\textbf{24.} ओरिन की पवन, \qty{31}{m/s}, \ang{45} N पर आंचलिक माध्य से अधिक
है और \ang{30} N पर उससे कम है। ध्रुवीय-अग्र जेट उत्तर और दक्षिण की ओर डोलता है;
किसी भी क्षण इसकी प्रवणता और इसकी पवन एक पट्टी पर संकेंद्रित होती है,
और अक्षांश वृत्त के चारों ओर औसत लेने से वह फैल जाती है।

\textbf{25.} $10.0 + 21.3$: स्टेपी के ऊपर \qty{10}{km} पर भू-विक्षेपी
पवन \textbf{\qty[text-series-to-math=true]{31}{m/s}} है, पश्चिम से
(गणना से \qty{31.3}{m/s})।
\end{solution}
